This is the central problem of this mini-project
\begin{problem}
	Let $p$ be an odd prime and let $k \equiv 2 \mod 4$. 
	We can write $k = r^2d$ where $d$ is square-free. 
	We want to investigate the integer solution of the following equation
\begin{equation*}
	y^p = x^2 + k
\end{equation*}
\end{problem}

\section{Q1}

In this section, we assume the class number of $\Q[\smd]$, denoted as $h(\Q[\smd])$, is not divisible by $p$, and we only work on coprime solutions $x, y$.

\subsection{General algorithm for coprime solutions when \texorpdfstring{$h(\Q[\smd])$ is not divisible by $p$}{}}

The diopantine equation can be factorised in $\Qd$ as 

$$
	y^p = (x+ r \smd) (x - r \smd)
$$

I claim that in $\O_{\Qd}$, $(x + r \smd)$ and $(x - r \smd)$ are coprime. 
Let $m$ be a common divisor of $(x + r \smd)$ and $(x - r \smd)$, then 
\begin{enumerate}
	\item $m$ is a divisor of their sum $2x$
	\item $m$ is a divisor of $y^p$ 
\end{enumerate}
If $m$ divides $x$, as $x$ and $y$ are coprime, $m$ must be a unit. 
If $m = 2$, then $y^p$ is even which implies $y$ is even. 
As $k$ is even, $x^2$ is even, which implies $x$ is even, contradicting the fact $x$ is coprime to $y$.
So $m$ can only be a unit, which proves the claim.

As a result, we can write $\langle x + r\smd \rangle$ as $\mathfrak{a}^p$ for some ideal $\mathfrak{a}$ in $\O_{\Qd}$. 

Since $h(\Q[\smd])$ is not divisible by $p$, $\a$ is a principal ideal, so we can find some $\alpha \in \O_{\Qd}$ such that $\langle \alpha\rangle = \a$. 
We can write $\alpha = a + b\smd$ for $a, b \in \Z$, then we have 
$(a + b\smd)^p = u(x + r\smd)$ for some unit $u$ in $\O_{\Qd}$. (We know the unit group of $\O_{\Qd}$ is $\{\pm 1\}$.)

Let us use the binomial expansion to expand $(a + b\smd)^p$ 
$$
	(a + b\smd)^p = a^p + \binom{p}{1} a^{p-1} b \smd + \binom{p}{2} a^{p-2} b^2 (\smd)^2 + \cdots + b^p (\smd)^p
$$
Note the term with $\smd$ can be written as 
\begin{dmath*}
	b\smd \left( \binom{p}{1} a^{p-1} + \binom{p}{3} a^{p-3} b^2 (\smd)^2 + \cdots + \binom{p}{p} b^{p-1} (\smd)^{p-1} \right) = b \smd f(a, b) = r
\end{dmath*}
where $f(a, b)$ is some homogeneous polynomial of degree $p-1$ in $a$ and $b$.
It is evident that $b$ is a divisor of $r$. 
We can try all possible values of $b$ and check if the corresponding $a$ is an integer.
In this way, we can find all possible integer solutions (if any exists) of the original diophantine equation.

\subsection{Example Solutions}

\textbf{Solution for} $y^3 = x^2 + 2$.

In this case $d = 2$ and the class number of $\Q[\sqrt{-2}]$ is 1, which is not divisible by 3.

By our procedure above, we can write the equation as
\begin{align*}
	y^3 = (x + \sqrt{-2})(x - \sqrt{-2})
\end{align*}
our argument shows that
$$
	x + \sqrt{-2} = (a + b\sqrt{-2})^3
$$
for some $a, b \in \Z$.
Expanding the right hand side, we have
\begin{align*}
	(a + b\sqrt{-2})^3 &= a^3 + 3a^2 b \sqrt{-2} - 6ab^2 - 2 \sqrt{-2} b^3 \\ 
					   &= (a^3 - 6ab^2) + b(3a^2 - 2b^2) \sqrt{-2}
\end{align*}
That is $b(3a^2 - 2b^2) = 1$ for some $a, b \in \Z$.
Since $b$ is a divisor of 1, it can be $\pm 1$.
\begin{enumerate}
	\item $b = 1$ gives $3a^2 - 2 = 1 \implies a = \pm 1$
	\item $b = -1$ gives $3a^2 - 2 = -1 \implies 3a^2 = 1$ which has no solution.
\end{enumerate}
So only options are $b = 1, a = \pm 1$.
Plugin to origin equation gives $y = (1 + \sqrt{-2})(1 - \sqrt{-2}) =  3$ and $x = \pm 5$.

\textbf{Solution for} $y^3 = x^2 + 18$.

In this case we have $p = 3, d = 2, r = 3$ and the equation can be factorised as
$$
	y^3 = (x + 3\sqrt{-2})(x - 3\sqrt{-2})
$$
By similar arguments as above we can write 
$$
	x + 3\sqrt{-2} = (a + b\sqrt{-2})^3
$$
which gives $b(3a^2 - 2b^2) = 3$, and $b$ has to be $\pm 1, \pm 3$ which are the only divisors of 3.
\begin{enumerate}
	\item $b = 1$ gives $3a^2 - 2 = 3 \implies 3a^2 = 5$ which has no solution.
	\item $b = -1$ gives $3a^2 - 2 = -3 \implies 3a^2 = -1$ which has no solution.
	\item $b = 3$ gives $3a^2 - 9 = 1 \implies 3a^2 = 10$ which has no solution.
	\item $b = -3$ gives $3a^2 - 9 = -1 \implies 3a^2 = 8$ which has no solution.
\end{enumerate}
So there is no solution for this equation.
